\begin{explanationbox}
  \textbf{\textcolor{CExplanation}{为什么会想到它？}}\par\nopagebreak[4]

  遇到 $f(a+t)=f(a-t)$，我们很容易看出两边的\textbf{函数值（高度）相等}。但题目往往追问的是 $f'$，也就是\textbf{图像在这些位置的斜率}。高度与斜率不是一回事，不过对称图像上的斜率有很整齐的规律。

  \medskip
  \textbf{\textcolor{CExplanation}{第一组：左右同高，左右斜率相反}}\par\nopagebreak[4]

  看 $f(x)=(x-2)^2+1$：$f(1)=f(3)=2$，但 $f'(1)=-2$，$f'(3)=2$。左边往右走是\textbf{下坡}，右边往右走是\textbf{上坡}，斜率大小一样、符号相反。
  \[
    \underbrace{f(a+t)=f(a-t)}_{\text{高度相同}}
    \quad\Longrightarrow\quad
    \underbrace{f'(a+t)=-f'(a-t)}_{\text{斜率反号}}.
  \]

  \medskip
  \textbf{\textcolor{CExplanation}{第二组：绕一个点对称，左右斜率相同}}\par\nopagebreak[4]

  看 $f(x)=(x-1)^3+7$：$f(0)=6$，$f(2)=8$，两个值以 $7$ 为中点；$f'(0)=f'(2)=3$。图中两条橙色切线平行，直接看出左右斜率相同。
  \[
    \underbrace{f(a+t)+f(a-t)=2f(a)}_{\text{左右函数值的平均数固定}}
    \quad\Longrightarrow\quad
    \underbrace{f'(a+t)=f'(a-t)}_{\text{斜率相同}}.
  \]

  \medskip
  \textbf{\textcolor{CExplanation}{图像对照：先看高度，再看斜率}}
  \begin{center}
  \begin{tikzpicture}[font=\small,>=stealth,line join=round]
\draw[draw=teal!45!black,fill=green!3,rounded corners=3pt] (0,0) rectangle (13.6,7.5);
\node at (6.8,6.95) {高度的对称性，怎样变成切线斜率的关系？};
\node at (3.4,6.35) {$f(x)=(x-2)^2+1$};
\node at (10.2,6.35) {$g(x)=(x-1)^3+7$};
\begin{scope}[shift={(1.2,1.25)},x=0.95cm,y=0.44cm]
\draw[gray,->] (-0.2,0)--(4.4,0) node[right] {$x$};
\draw[gray,->] (0,-0.3)--(0,10.5);
\draw[teal!75!black,very thick,domain=0.3:3.7,samples=100] plot(\x,{(\x-2)^2+1});
\draw[orange!85!black,thick,domain=.55:1.45] plot(\x,{4-2*\x});
\draw[orange!85!black,thick,domain=2.55:3.45] plot(\x,{2*\x-4});
\draw[dashed,gray] (2,0)--(2,8.8);
\fill[teal!75!black] (1,2) circle(2.5pt);
\fill[teal!75!black] (3,2) circle(2.5pt);
\node at (.6,6.4) {斜率 $-2$};\node at (3.5,6.4) {斜率 $2$};
\node[below] at (2,-.1) {$2$};
\end{scope}
\begin{scope}[shift={(8.5,1.25)},x=1cm,y=0.44cm]
\draw[gray,->] (-0.3,0)--(3.9,0) node[right] {$x$};
\draw[gray,->] (0,-0.3)--(0,10.5);
\draw[teal!75!black,very thick,domain=-0.2:2.3,samples=100] plot(\x,{(\x-1)^3+7});
\draw[orange!85!black,thick,domain=-.25:.35] plot(\x,{3*\x+6});
\draw[orange!85!black,thick,domain=1.65:2.25] plot(\x,{3*\x+2});
\draw[gray,dotted] (0,6)--(2,8);
\fill[teal!75!black] (0,6) circle(2.5pt);
\fill[teal!75!black] (2,8) circle(2.5pt);
\fill[blue!70!black] (1,7) circle(2.5pt);
\node[left] at (-.1,6) {$(0,6)$};\node[right] at (2.15,8) {$(2,8)$};
\node[below right] at (1.1,6.2) {$(1,7)$};
\node at (0,10.2) {斜率 $3$};\node at (2.6,10.2) {斜率 $3$};
\draw[gray] (0,0+.08)--(0,0-.08);\node[below] at (0,0-.15) {$0$};
\draw[gray] (1,0+.08)--(1,0-.08);\node[below] at (1,0-.15) {$1$};
\draw[gray] (2,0+.08)--(2,0-.08);\node[below] at (2,0-.15) {$2$};
\end{scope}
\node at (3.4,.45) {轴对称：左右切线斜率反号};
\node at (10.2,.45) {中心对称：左右切线斜率相同};
\end{tikzpicture}

  \end{center}

  \medskip
  \textbf{\textcolor{CExplanation}{为什么可以反推？}}\par\nopagebreak[4]

  若两边斜率一直反号，则 $f(a+t)-f(a-t)$ 随 $t$ 变化时导数为零，所以差是常数；代入 $t=0$ 可知常数是 $0$，于是两边同高。若两边斜率一直相同，则 $f(a+t)+f(a-t)$ 是常数；代入 $t=0$，这个常数就是 $2f(a)$。

  \medskip
  \textbf{\textcolor{CExplanation}{容易忽略的推广：中间不一定是 0}}\par\nopagebreak[4]

  导函数关于 $(a,b)$ 中心对称，表示左右斜率的\textbf{平均数是 $b$}：
  \[
    \frac{f'(a+t)+f'(a-t)}{2}=b.
  \]
  因此原函数左右函数值的差是 $2bt$。$b=0$ 才能保证关于 $x=a$ 轴对称；$b\ne0$ 只能排除\textbf{这条轴}，不能排除所有轴。

  \medskip
  \textbf{\textcolor{CExplanation}{考试中什么时候最快？}}
  \begin{itemize}[leftmargin=*,itemsep=0.35em,topsep=0.3em]
    \item \textbf{只给对称条件：} 直接写 $a\pm t$ 的等式，转为导函数条件，不必求 $f(x)$ 的解析式。
    \item \textbf{只给导函数的对称轴：} 立刻知道原函数的中心横坐标是 $a$，纵坐标要写 $f(a)$。
    \item \textbf{同时给两类对称：} 把对称式连续代入，可能得到 $f(x+T)=f(x)$ 等周期关系。
  \end{itemize}
\end{explanationbox}
